% =====================================================================
% CM Séance 2 — Analyse algorithmique
% Structures de Données et Algorithmes Avancés — M1 — ESP/UCAD
% Dr. El Hadji Bassirou TOURÉ — 2025–2026
% Compilation : pdflatex (2 passes)
% =====================================================================
\documentclass[11pt,a4paper]{article}

\usepackage[utf8]{inputenc}
\usepackage[T1]{fontenc}
\usepackage{lmodern}
\usepackage[french]{babel}

\usepackage[margin=2.5cm]{geometry}
\usepackage{amsmath,amssymb,amsthm}
\usepackage{graphicx}
\graphicspath{{figures/}}
\usepackage{xcolor}
\usepackage{booktabs}
\usepackage{array}
\usepackage{enumitem}
\usepackage{fancyhdr}
\usepackage{longtable}
\usepackage{multirow}
\usepackage{pifont}
\usepackage{listings}
\usepackage{tcolorbox}
\tcbuselibrary{theorems,breakable,skins}
\usepackage{hyperref}

\hypersetup{colorlinks=true,linkcolor=blue!60!black,urlcolor=blue!60!black,
            citecolor=blue!60!black}

\newcommand{\cmark}{\ding{51}}
\newcommand{\xmark}{\ding{55}}
\newcommand{\bigO}{\mathcal{O}}
\newcommand{\bigOmega}{\Omega}
\newcommand{\bigTheta}{\Theta}

% ===== Boîtes =====
\newtcolorbox{objectifbox}[1][]{colback=blue!5,colframe=blue!50,
    fonttitle=\bfseries,title=#1,breakable}
\newtcolorbox{infobox}[1][]{colback=green!5,colframe=green!50!black,
    fonttitle=\bfseries,title=#1,breakable}
\newtcolorbox{attentionbox}[1][]{colback=orange!5,colframe=orange!70,
    fonttitle=\bfseries,title=#1,breakable}
\newtcolorbox{retenirbox}[1][]{colback=yellow!10,colframe=yellow!50!black,
    fonttitle=\bfseries,title=#1,breakable}
\newtcolorbox{prereqbox}[1][]{colback=cyan!5,colframe=cyan!50!black,
    fonttitle=\bfseries,title=#1,breakable}
\newtcolorbox{theorembox}[1][]{colback=purple!5,colframe=purple!50,
    fonttitle=\bfseries,title=#1,breakable}
\newtcolorbox{appliboxQ}[1][]{colback=gray!8,colframe=gray!60,
    fonttitle=\bfseries,title=#1,breakable}
\newtcolorbox{appliboxR}[1][]{colback=green!3,colframe=green!40!black,
    fonttitle=\bfseries\itshape,title=#1,breakable}
\newtcolorbox{exercicebox}[1][]{colback=white,colframe=gray!70,
    fonttitle=\bfseries,title=#1,breakable}
\newtcolorbox{lienbox}[1][]{colback=blue!3,colframe=blue!70!black,
    fonttitle=\bfseries\itshape,title=#1,breakable,
    before upper={\itshape}}

% ===== Théorèmes =====
\theoremstyle{definition}
\newtheorem{definition}{Définition}[section]
\newtheorem{theorem}{Théorème}[section]
\newtheorem{proposition}[theorem]{Proposition}
\theoremstyle{remark}
\newtheorem{example}{Exemple}[section]
\newtheorem{remark}{Remarque}[section]

% ===== Code =====
\lstset{
  language=Java,
  basicstyle=\ttfamily\small,
  backgroundcolor=\color{gray!8},
  frame=single,
  framerule=0.5pt,
  rulecolor=\color{gray!40},
  keywordstyle=\color{blue!70!black}\bfseries,
  commentstyle=\color{green!50!black}\itshape,
  stringstyle=\color{red!60!black},
  numbers=left,
  numberstyle=\tiny\color{gray},
  breaklines=true,
  breakatwhitespace=true,
  showstringspaces=false,
  tabsize=4,
  extendedchars=true,
  xleftmargin=4pt,
  xrightmargin=4pt,
  aboveskip=8pt,
  belowskip=8pt,
  literate={é}{{\'e}}1 {è}{{\`e}}1 {ê}{{\^e}}1 {à}{{\`a}}1 {ù}{{\`u}}1
           {ô}{{\^o}}1 {î}{{\^i}}1 {ç}{{\c{c}}}1 {û}{{\^u}}1 {â}{{\^a}}1
           {ë}{{\"e}}1 {ï}{{\"i}}1 {É}{{\'E}}1 {È}{{\`E}}1 {À}{{\`A}}1
           {Ç}{{\c{C}}}1 {Ê}{{\^E}}1 {œ}{{\oe}}1 {°}{{$^\circ$}}1
           {«}{{\guillemotleft}}1 {»}{{\guillemotright}}1
}

% ===== En-têtes =====
\pagestyle{fancy}
\fancyhf{}
\fancyhead[L]{\small S02 --- Analyse algorithmique}
\fancyhead[R]{\small Structures de Données Avancées --- M1}
\fancyfoot[L]{\small Dr. El Hadji Bassirou TOURÉ}
\fancyfoot[C]{\thepage}
\fancyfoot[R]{\small 2025--2026}
\renewcommand{\headrulewidth}{0.4pt}
\renewcommand{\footrulewidth}{0.4pt}

\title{%
\vspace{-1cm}
\textbf{\LARGE Structures de Données et}\\
\textbf{\LARGE Algorithmes Avancés}\\[0.8cm]
\textbf{\huge Séance 2 --- Analyse algorithmique}\\[0.3cm]
\textbf{\Large Comment mesurer l'efficacité d'un algorithme ?}\\[0.5cm]
\large Master 1 --- 2025--2026}
\author{Dr. El Hadji Bassirou TOURÉ\\
Département Génie Informatique\\
Université Cheikh Anta Diop de Dakar}
\date{}

\begin{document}
\maketitle
\thispagestyle{empty}

\begin{prereqbox}[Navigation conceptuelle]
\textbf{Ce que vous savez déjà :} séparer un contrat de son implémentation
(S1), et réaliser List, Stack, Queue et Deque par tableau ou par nœuds
chaînés.\\[4pt]
\textbf{La limite qu'on dépasse aujourd'hui :} en S1, nous avons répété que
\texttt{remove(0)} est « coûteux » et que \texttt{get(i)} est « constant »
sans rien démontrer. Une intuition n'est pas un argument, et deux ingénieurs
qui s'appuient sur leurs intuitions ne peuvent pas trancher.\\[4pt]
\textbf{Ce que vous saurez faire à la fin :}
\begin{enumerate}[nosep]
    \item Établir la fonction de coût $f(n)$ d'un fragment de code par
          comptage granulaire des opérations élémentaires.
    \item Énoncer les définitions de $\bigO$, $\bigOmega$, $\bigTheta$ et
          démontrer une borne en exhibant $c$ et $n_0$.
    \item Distinguer meilleur cas, pire cas et cas moyen, et justifier
          pourquoi le pire cas est la référence.
    \item Résoudre une récurrence de diviser-pour-régner par le théorème
          principal.
\end{enumerate}
\textbf{Ouvre la voie vers :} S3, où l'on cherchera une structure permettant
l'accès par clé sans parcourir les données.
\end{prereqbox}

\tableofcontents
\newpage

% =====================================================================
\section{D'où l'on part}
% =====================================================================

À la fin de la séance précédente, un constat nous est resté sur les bras :
une file construite sur une \texttt{ArrayList} avec \texttt{remove(0)} est
parfaitement \emph{correcte} --- elle sert les clients dans le bon ordre ---
et pourtant nous la déclarons mauvaise. Sur quoi repose ce jugement ? Pas sur
les tests, qui passent. Pas sur le chronomètre, qui dépend de la machine, du
système et de l'humeur du processeur.

Il nous faut donc une grandeur qui ne dépende ni du matériel ni du langage,
et qui dise comment le coût d'un algorithme \emph{évolue} quand la taille des
données augmente. Cette grandeur s'obtient en deux temps : on compte les
opérations élémentaires, ce qui donne une fonction $f(n)$ ; puis on classe
cette fonction dans une famille de croissance. Le premier temps est
arithmétique, le second est mathématique. Les deux ensemble forment l'outil
qui nous servira jusqu'à la fin du semestre.

% =====================================================================
\section{Compter les opérations}
% =====================================================================

\subsection{La convention du cours}

Compter suppose de dire ce qu'on compte. Nous adoptons une convention
\emph{granulaire} : chaque opération élémentaire compte pour une unité, y
compris les affectations et les accès mémoire.

\begin{table}[h!]
\centering
\caption{Convention de comptage retenue dans ce cours}
\begin{tabular}{ll}
\toprule
\textbf{Instruction} & \textbf{Opérations comptées} \\
\midrule
\texttt{s = s + i}    & 2 : une addition, une affectation \\
\texttt{i++}          & 2 : une addition, une affectation \\
\texttt{tab[i]}       & 1 : un accès mémoire \\
\texttt{i < n}        & 1 : une comparaison \\
\texttt{s = 0}        & 1 : une affectation \\
\texttt{return x}     & 1 \\
\bottomrule
\end{tabular}
\end{table}

\begin{attentionbox}[Une convention explicite plutôt qu'une vérité]
D'autres cours comptent \texttt{s = s + i} pour une seule opération. Le
résultat final --- la classe de complexité --- est le même, car les
constantes disparaissent à l'étape suivante. Mais tant qu'on manipule
$f(n)$, la convention doit être explicite : un corrigé ne peut pas être
discuté si la règle du jeu n'est pas écrite. \textbf{Dans ce cours, on
compte comme dans le tableau ci-dessus, et on l'annonce.}
\end{attentionbox}

\subsection{Premier exemple : une boucle simple}

\begin{lstlisting}[caption={Somme des $n$ premiers entiers}]
int s = 0;                      // 1 affectation
for (int i = 0; i < n; i++) {   // 1 affectation + (n+1) comparaisons + n*(add+aff)
    s = s + i;                  // n * (1 addition + 1 affectation)
}
\end{lstlisting}

Le décompte, ligne par ligne :
\begin{align*}
\texttt{s = 0} &: 1 \\
\texttt{i = 0} &: 1 \\
\texttt{i < n} &: n+1 \quad \text{($n$ fois vraie, une fois fausse)} \\
\texttt{s = s + i} &: 2n \\
\texttt{i++} &: 2n
\end{align*}
\[
f(n) \;=\; 1 + 1 + (n+1) + 2n + 2n \;=\; 5n + 3
\]

\begin{figure}[h!]
\centering
\includegraphics[width=0.88\textwidth]{fig_01.pdf}
\caption{La courbe bleue provient d'un compteur incrémenté à chaque
opération lors d'une exécution réelle ; la droite orange est la formule
$5n+3$. Les deux coïncident pour $n = 0, \dots, 10$.}
\end{figure}

\begin{infobox}[Pourquoi vérifier une formule de comptage]
Une formule de comptage se trompe facilement d'une constante : il est très
courant d'oublier la comparaison finale qui fait sortir de la boucle, ou de
compter $n$ au lieu de $n+1$. Instrumenter le code --- un compteur
incrémenté à chaque opération --- coûte quelques lignes et tranche la
question. C'est ce que vous ferez au Lab.
\end{infobox}

\begin{appliboxQ}[Mini-exercice 1 --- La comparaison oubliée]
Un étudiant obtient $f(n) = 5n + 2$. Quelle opération a-t-il omise ?
\end{appliboxQ}

\begin{appliboxR}[Réponse]
La dernière évaluation de \texttt{i < n}, celle qui est fausse et qui fait
sortir de la boucle : il a compté $n$ comparaisons au lieu de $n+1$.

\textbf{Le piège classique :} raisonner « une comparaison par tour de
boucle ». Il y a toujours une comparaison de plus que de tours.
\end{appliboxR}

\subsection{Boucles imbriquées}

\begin{lstlisting}[caption={Boucle triangulaire}]
int c = 0;
for (int i = 0; i < n; i++) {
    for (int j = i; j < n; j++) {
        c = c + 1;
    }
}
\end{lstlisting}

La boucle interne s'exécute $n - i$ fois pour chaque $i$, soit au total
\[
\sum_{i=0}^{n-1} (n-i) \;=\; n + (n-1) + \dots + 1 \;=\; \frac{n(n+1)}{2}
\]
tours. En appliquant la convention, on obtient
\[
f(n) \;=\; 2{,}5\,n^2 + 7{,}5\,n + 3
\]
Pour $n = 5$ : $f(5) = 62{,}5 + 37{,}5 + 3 = 103$, ce que confirme
l'instrumentation du code.

\begin{appliboxQ}[Mini-exercice 2 --- Le terme dominant]
Dans $f(n) = 2{,}5n^2 + 7{,}5n + 3$, quelle est la part du terme en $n^2$
pour $n = 10$ ? Pour $n = 1000$ ?
\end{appliboxQ}

\begin{appliboxR}[Réponse]
$n = 10$ : $250$ sur $328$, soit environ $76\,\%$.\\
$n = 1000$ : $2\,500\,000$ sur $2\,507\,503$, soit environ $99{,}7\,\%$.

\textbf{La leçon :} plus $n$ grandit, plus le terme de plus haut degré
écrase les autres. C'est exactement ce constat que la notation $\bigO$
formalise --- et c'est aussi pourquoi il est inutile de s'écharper sur les
constantes.
\end{appliboxR}

\subsection{Une boucle qui ne progresse pas d'un pas}

\begin{lstlisting}[caption={Progression multiplicative}]
for (int i = 1; i < n; i = i * 3) {
    traiter(i);
}
\end{lstlisting}

Les valeurs prises par \texttt{i} sont $1, 3, 9, 27, 81, \dots$, c'est-à-dire
$3^0, 3^1, 3^2, \dots$ La boucle s'arrête dès que $3^k \geq n$, donc après
environ $\log_3 n$ tours. Pour $n = 100$ : cinq tours ($1, 3, 9, 27, 81$).
Pour $n = 10\,000$ : neuf tours.

\begin{attentionbox}[L'erreur la plus fréquente de la séance]
Lire \texttt{i = i * 3} comme une boucle linéaire. Multiplier $n$ par $10$
n'ajoute ici que deux tours, alors qu'une boucle en \texttt{i++} en
ajouterait $9n$. Test de reconnaissance : \emph{la variable de boucle
est-elle additionnée ou multipliée ?}
\end{attentionbox}

\subsection{Le coût qu'on ne voit pas dans le code}

\begin{lstlisting}[caption={Deux lignes, deux coûts très différents}]
for (int i = 0; i < n; i++) {
    liste.add(0, i);   // insertion en tête
}
\end{lstlisting}

La ligne ne comporte aucune boucle visible. Pourtant, si \texttt{liste} est
une \texttt{ArrayList}, l'insertion en position $0$ décale les $i$ éléments
déjà présents : le $i$-ième tour effectue $i$ déplacements, et le total vaut
$0 + 1 + \dots + (n-1) = n(n-1)/2$.

\begin{figure}[h!]
\centering
\includegraphics[width=0.88\textwidth]{fig_05.pdf}
\caption{Coût cumulé de $n$ insertions en tête, calculé en simulant les
déplacements. Pour $n = 2000$ : $1\,999\,000$ déplacements contre $2000$
rebranchements de références.}
\end{figure}

\begin{appliboxQ}[Mini-exercice 3 --- Le facteur en jeu]
Pour $n = 2000$, quel est le rapport entre les deux courbes de la figure ?
Et pour $n = 20\,000$ ?
\end{appliboxQ}

\begin{appliboxR}[Réponse]
$n = 2000$ : $1\,999\,000 / 2000 \approx 1000$.\\
$n = 20\,000$ : $199\,990\,000 / 20\,000 \approx 10\,000$.

\textbf{La leçon :} l'écart n'est pas un facteur fixe, il grandit avec $n$.
Un programme qui « marche bien en démonstration » sur $100$ éléments peut
être inutilisable sur $100\,000$. C'est pourquoi on analyse au lieu de
tester une seule taille.
\end{appliboxR}

% =====================================================================
\section{Les notations asymptotiques}
% =====================================================================

\subsection{Big-O : une borne supérieure}

\begin{definition}[Notation $\bigO$]
Soient $f, g : \mathbb{N} \to \mathbb{R}^{+}$. On écrit
$f(n) \in \bigO(g(n))$ s'il existe deux constantes $c > 0$ et
$n_0 \in \mathbb{N}$ telles que
\[
\forall n \geq n_0, \quad f(n) \leq c \cdot g(n).
\]
\end{definition}

Trois éléments méritent d'être soulignés. La constante $c$ autorise à
ignorer les facteurs multiplicatifs. Le seuil $n_0$ autorise à ignorer le
comportement sur les petites tailles. Et l'appartenance est une
\emph{majoration} : dire $f \in \bigO(n^2)$ n'interdit pas que $f$ soit en
réalité linéaire.

\begin{figure}[h!]
\centering
\includegraphics[width=0.85\textwidth]{fig_02.pdf}
\caption{Au-delà de $n_0 = 13$, la courbe $5n^2$ reste au-dessus de
$3n^2 + 20n + 50$. Le seuil est calculé, pas estimé : en $n = 12$ la
majoration est fausse ($722 > 720$), en $n = 13$ elle devient vraie
($817 \leq 845$).}
\end{figure}

\begin{example}[Une preuve rédigée]
Montrons que $f(n) = 3n^2 + 20n + 50 \in \bigO(n^2)$.

Pour tout $n \geq 1$, on a $n \leq n^2$ et $1 \leq n^2$, donc
\[
3n^2 + 20n + 50 \;\leq\; 3n^2 + 20n^2 + 50n^2 \;=\; 73\,n^2.
\]
Avec $c = 73$ et $n_0 = 1$, la définition est satisfaite. \hfill $\square$
\end{example}

\begin{infobox}[Le couple $c$ et $n_0$ n'est pas unique]
La preuve ci-dessus donne $(c, n_0) = (73, 1)$ ; la figure exhibe
$(5, 13)$. Les deux sont corrects : il suffit d'exhiber \emph{un} couple.
Chercher le plus petit $c$ possible est un exercice d'analyse, pas une
exigence de la définition.
\end{infobox}

\begin{appliboxQ}[Mini-exercice 4 --- Une preuve à produire]
Montrez que $f(n) = 7n + 40 \in \bigO(n)$ en exhibant un couple
$(c, n_0)$.
\end{appliboxQ}

\begin{appliboxR}[Réponse]
Pour $n \geq 1$, $40 \leq 40n$, donc $7n + 40 \leq 47n$ : le couple
$(47, 1)$ convient. On peut faire mieux : pour $n \geq 40$, $40 \leq n$,
donc $7n + 40 \leq 8n$, et $(8, 40)$ convient aussi.

\textbf{Remarque :} plus on repousse $n_0$, plus on peut serrer $c$. C'est
tout le sens du mot « asymptotique ».
\end{appliboxR}

\subsection{$\bigOmega$ et $\bigTheta$}

\begin{definition}[Notations $\bigOmega$ et $\bigTheta$]
$f(n) \in \bigOmega(g(n))$ s'il existe $c > 0$ et $n_0$ tels que
$f(n) \geq c\,g(n)$ pour tout $n \geq n_0$.\\
$f(n) \in \bigTheta(g(n))$ si $f(n) \in \bigO(g(n))$ \emph{et}
$f(n) \in \bigOmega(g(n))$.
\end{definition}

$\bigO$ majore, $\bigOmega$ minore, $\bigTheta$ encadre. Quand on connaît le
coût exact d'un algorithme, c'est $\bigTheta$ qu'il faut écrire ; $\bigO$
reste correct mais dit moins.

\begin{proposition}[Propriétés utiles]
Pour $f, g, h$ à valeurs positives :
\begin{enumerate}[nosep,label=(\roman*)]
    \item \textbf{Transitivité :} si $f \in \bigO(g)$ et $g \in \bigO(h)$,
          alors $f \in \bigO(h)$.
    \item \textbf{Somme :} $\bigO(f) + \bigO(g) = \bigO(\max(f, g))$.
    \item \textbf{Produit :} si $f_1 \in \bigO(g_1)$ et
          $f_2 \in \bigO(g_2)$, alors $f_1 f_2 \in \bigO(g_1 g_2)$.
    \item \textbf{Constantes :} $\bigO(c \cdot f) = \bigO(f)$ pour toute
          constante $c > 0$.
\end{enumerate}
\end{proposition}

La propriété (ii) justifie la pratique courante : dans une suite
d'instructions, seule la plus coûteuse compte. La propriété (iii) justifie
l'analyse des boucles imbriquées par multiplication des coûts.

\begin{figure}[h!]
\centering
\includegraphics[width=\textwidth]{fig_03.pdf}
\caption{Les classes usuelles. À $n = 100$ : $\log_2 n \approx 6{,}6$,
$n\log_2 n \approx 664$, $n^2 = 10\,000$. L'échelle logarithmique, à droite,
rend les cinq familles lisibles simultanément.}
\end{figure}

\begin{appliboxQ}[Mini-exercice 5 --- Vrai ou faux]
\begin{enumerate}[nosep,label=(\alph*)]
    \item $5n^2 + 3n \in \bigO(n^3)$
    \item $5n^2 + 3n \in \bigTheta(n^3)$
    \item $n^2 \in \bigOmega(n \log n)$
    \item $2^{n+1} \in \bigO(2^n)$
\end{enumerate}
\end{appliboxQ}

\begin{appliboxR}[Réponse]
(a) vrai --- $\bigO$ est une majoration, elle n'a pas à être serrée.\\
(b) faux --- il manque la minoration : $n^2$ ne croît pas comme $n^3$.\\
(c) vrai --- $n^2$ domine $n\log n$.\\
(d) vrai --- $2^{n+1} = 2 \cdot 2^n$, et une constante multiplicative
disparaît.

\textbf{Le piège classique :} confondre (a) et (b). Une majoration vraie
mais lâche n'est pas fausse ; elle est simplement peu informative. La
question d'examen porte souvent sur cette nuance.
\end{appliboxR}

\subsection{Du comptage à la classe}

La démarche complète tient en trois gestes : compter, garder le terme
dominant, oublier les constantes.
\[
f(n) = 2{,}5n^2 + 7{,}5n + 3
\;\longrightarrow\;
\text{terme dominant } 2{,}5n^2
\;\longrightarrow\;
f(n) \in \bigTheta(n^2)
\]

\begin{appliboxQ}[Mini-exercice 6 --- Classer sans calculer]
Donnez la classe $\bigTheta$ de chaque fragment :
\begin{enumerate}[nosep,label=(\alph*)]
    \item une boucle \texttt{for (i = 0; i < n; i++)} contenant une
          instruction de coût constant ;
    \item deux boucles successives sur $n$ ;
    \item une boucle sur $n$ contenant une boucle sur $n$ ;
    \item une boucle \texttt{for (i = 1; i < n; i = i * 2)}.
\end{enumerate}
\end{appliboxQ}

\begin{appliboxR}[Réponse]
(a) $\bigTheta(n)$ --- (b) $\bigTheta(n)$ : deux boucles successives
s'additionnent, $n + n = 2n$ --- (c) $\bigTheta(n^2)$ : les boucles
imbriquées se multiplient --- (d) $\bigTheta(\log n)$.

\textbf{La règle :} successives $\rightarrow$ on additionne (et le maximum
l'emporte) ; imbriquées $\rightarrow$ on multiplie.
\end{appliboxR}

% =====================================================================
\section{Analyse de cas}
% =====================================================================

Pour une même taille $n$, le coût dépend souvent du contenu des données.

\begin{lstlisting}[caption={Recherche linéaire}]
public static int chercher(int[] tab, int cible) {
    for (int i = 0; i < tab.length; i++) {
        if (tab[i] == cible) return i;
    }
    return -1;
}
\end{lstlisting}

\begin{table}[h!]
\centering
\caption{Les trois cas de la recherche linéaire}
\begin{tabular}{lll}
\toprule
\textbf{Cas} & \textbf{Situation} & \textbf{Coût} \\
\midrule
Meilleur & la cible est en case $0$ & $\bigTheta(1)$ \\
Pire & la cible est en dernière case ou absente & $\bigTheta(n)$ \\
Moyen & cible présente, position uniforme & $\approx n/2$ comparaisons, soit $\bigTheta(n)$ \\
\bottomrule
\end{tabular}
\end{table}

Le pire cas sert de référence pour deux raisons. Il fournit une garantie
--- « jamais plus que cela » --- sur laquelle on peut dimensionner un
système. Et il ne suppose rien sur la distribution des données, alors que le
cas moyen exige une hypothèse probabiliste souvent invérifiable en pratique.

\begin{attentionbox}[Cas et notation sont deux questions distinctes]
« Pire cas » désigne une famille d'entrées ; $\bigO$ désigne une famille de
fonctions. On peut donc parler du $\bigTheta$ du meilleur cas comme du
$\bigO$ du pire cas. Écrire « $\bigO$ signifie pire cas » est une confusion
fréquente, et fausse.
\end{attentionbox}

\begin{appliboxQ}[Mini-exercice 7 --- Identifier le pire cas]
Pour chacun, décrivez l'entrée qui réalise le pire cas :
\begin{enumerate}[nosep,label=(\alph*)]
    \item recherche d'un doublon par double boucle ;
    \item insertion de $n$ clés dans un tableau trié par insertion directe.
\end{enumerate}
\end{appliboxQ}

\begin{appliboxR}[Réponse]
(a) un tableau sans aucun doublon : aucune sortie anticipée, les deux
boucles vont jusqu'au bout.\\
(b) des clés fournies en ordre décroissant : chaque nouvelle clé doit être
insérée en tête, donc décale tout.

\textbf{Remarque :} le pire cas est souvent le cas « ordonné à l'envers »
ou le cas « absent ». Ce sont les entrées à mettre dans vos tests de
performance.
\end{appliboxR}

% =====================================================================
\section{Analyser le code récursif}
% =====================================================================

\subsection{Poser la récurrence}

Le coût d'une fonction récursive s'exprime en fonction de lui-même.

\begin{lstlisting}[caption={Recherche dichotomique}]
static int dicho(int[] tab, int cible, int gauche, int droite) {
    if (gauche > droite) return -1;
    int milieu = (gauche + droite) / 2;
    if (tab[milieu] == cible) return milieu;
    if (tab[milieu] < cible) return dicho(tab, cible, milieu + 1, droite);
    else                     return dicho(tab, cible, gauche, milieu - 1);
}
\end{lstlisting}

Chaque appel réduit l'intervalle de moitié et effectue un travail constant :
\[
T(n) = T(n/2) + \bigTheta(1), \qquad T(1) = \bigTheta(1).
\]
En déroulant : après $k$ appels, l'intervalle a pour taille $n/2^k$, et la
récursion s'arrête quand $n/2^k = 1$, soit $k = \log_2 n$. D'où
$T(n) = \bigTheta(\log n)$.

\subsection{Le théorème principal}

\begin{theorembox}[Théorème principal (Master Theorem)]
Soit $T(n) = a\,T(n/b) + f(n)$ avec $a \geq 1$, $b > 1$, et posons
$\alpha = \log_b a$.
\begin{enumerate}[nosep,label=\textbf{Cas \arabic*.}]
    \item Si $f(n) \in \bigO(n^{\alpha - \varepsilon})$ pour un
          $\varepsilon > 0$, alors $T(n) \in \bigTheta(n^{\alpha})$.
    \item Si $f(n) \in \bigTheta(n^{\alpha})$, alors
          $T(n) \in \bigTheta(n^{\alpha} \log n)$.
    \item Si $f(n) \in \bigOmega(n^{\alpha + \varepsilon})$ pour un
          $\varepsilon > 0$, et si $a\,f(n/b) \leq c\,f(n)$ pour un
          $c < 1$ et $n$ assez grand, alors $T(n) \in \bigTheta(f(n))$.
\end{enumerate}
\end{theorembox}

L'énoncé se lit comme une comparaison entre deux forces : le travail des
appels récursifs, mesuré par $n^{\alpha}$, et le travail effectué à chaque
niveau, mesuré par $f(n)$. Si les feuilles l'emportent, on est au cas 1 ; si
l'équilibre règne, au cas 2, et un facteur $\log n$ apparaît ; si la racine
l'emporte, au cas 3.

\begin{figure}[h!]
\centering
\includegraphics[width=\textwidth]{fig_04.pdf}
\caption{Coût par niveau de récursion pour $n = 1024$. À gauche le travail
se concentre sur les feuilles (total $2047$) ; au centre chaque niveau coûte
$1024$, pour $11$ niveaux (total $11\,264$) ; à droite la racine absorbe
presque tout (total $2\,096\,128$).}
\end{figure}

\begin{example}[Mergesort]
Trier par fusion coupe le tableau en deux et fusionne en temps linéaire :
$T(n) = 2T(n/2) + \bigTheta(n)$. Ici $a = 2$, $b = 2$, donc
$\alpha = \log_2 2 = 1$ et $f(n) = \bigTheta(n) = \bigTheta(n^{\alpha})$ :
c'est le cas 2, et $T(n) \in \bigTheta(n \log n)$. Nous reviendrons sur cet
algorithme en S6.
\end{example}

\begin{appliboxQ}[Mini-exercice 8 --- Appliquer le théorème]
Donnez la classe de $T(n)$ dans chaque cas :
\begin{enumerate}[nosep,label=(\alph*)]
    \item $T(n) = 2T(n/2) + \bigTheta(1)$
    \item $T(n) = T(n/2) + \bigTheta(1)$
    \item $T(n) = 4T(n/2) + \bigTheta(n)$
\end{enumerate}
\end{appliboxQ}

\begin{appliboxR}[Réponse]
(a) $\alpha = 1$, $f(n) = \bigTheta(1) = \bigO(n^{1-\varepsilon})$ : cas 1,
donc $\bigTheta(n)$.\\
(b) $a = 1$, $\alpha = \log_2 1 = 0$, $f(n) = \bigTheta(n^0)$ : cas 2, donc
$\bigTheta(\log n)$ --- c'est la dichotomie.\\
(c) $\alpha = \log_2 4 = 2$, $f(n) = \bigTheta(n) = \bigO(n^{2-\varepsilon})$ :
cas 1, donc $\bigTheta(n^2)$.

\textbf{Le piège classique :} confondre (a) et (b). Le nombre d'appels
récursifs $a$ change tout : un seul appel donne $\bigTheta(\log n)$, deux
appels donnent $\bigTheta(n)$.
\end{appliboxR}

\begin{appliboxQ}[Mini-exercice 9 --- Revenir à la question de départ]
Reprenez la file construite sur une \texttt{ArrayList} avec
\texttt{remove(0)} (S1). Établissez le coût total de $n$ opérations
\texttt{dequeue}, puis comparez à la file sur tableau circulaire.
\end{appliboxQ}

\begin{appliboxR}[Réponse]
Le $k$-ième \texttt{remove(0)} décale les éléments restants : le total est de
l'ordre de $n^2/2$ déplacements, soit $\bigTheta(n^2)$ pour la séquence
complète. Avec le tableau circulaire, chaque \texttt{dequeue} déplace un
indice : total $\bigTheta(n)$.

\textbf{La conclusion de la séance :} l'intuition de S1 est maintenant un
énoncé démontré, et le gain est quantifié --- d'un facteur proportionnel à
$n$, donc arbitrairement grand.
\end{appliboxR}

% =====================================================================
\begin{retenirbox}[Points clés à retenir]
\begin{enumerate}[nosep]
    \item Analyser, c'est d'abord compter. La convention de comptage doit
          être annoncée : ici \texttt{s = s + i} vaut deux opérations, et il
          y a toujours une comparaison de plus que de tours de boucle.
    \item $f(n) \in \bigO(g(n))$ signifie : il existe $c > 0$ et $n_0$ tels
          que $f(n) \leq c\,g(n)$ pour tout $n \geq n_0$. Prouver, c'est
          exhiber un couple $(c, n_0)$ --- pas nécessairement le meilleur.
    \item $\bigO$ majore, $\bigOmega$ minore, $\bigTheta$ encadre. Quand le
          coût est connu exactement, on écrit $\bigTheta$.
    \item Boucles successives : on additionne, le maximum l'emporte.
          Boucles imbriquées : on multiplie. Variable multipliée dans le
          pas : c'est logarithmique.
    \item Le pire cas est la référence, parce qu'il donne une garantie et
          ne suppose rien sur les données. « Pire cas » et « $\bigO$ » sont
          deux notions distinctes.
    \item Le code récursif se traduit en récurrence $T(n) = aT(n/b) + f(n)$,
          que le théorème principal résout en comparant $f(n)$ à
          $n^{\log_b a}$.
    \item Un coût peut être invisible dans le code : \texttt{liste.add(0, x)}
          tient en une ligne et cache un décalage complet.
\end{enumerate}
\end{retenirbox}

% =====================================================================
\section*{Ce qui vient ensuite : Tables de hachage}
\addcontentsline{toc}{section}{Ce qui vient ensuite : Tables de hachage}
% =====================================================================

\begin{prereqbox}[Vers la séance suivante]
\textbf{Ce qu'on sait faire maintenant :} établir $f(n)$, prouver une borne
asymptotique, distinguer les cas, et résoudre une récurrence de
diviser-pour-régner.\\[4pt]
\textbf{La limite qu'on va dépasser :} nos outils confirment une mauvaise
nouvelle. Dans un tableau, l'accès \emph{par rang} est en $\bigTheta(1)$,
mais la recherche \emph{par valeur} est en $\bigTheta(n)$ ; les nœuds
chaînés ne font pas mieux. Or la plupart des applications cherchent par clé
--- un numéro d'étudiant, une plaque de bus --- et non par position.\\[4pt]
\textbf{En S3 :} on construit une structure qui calcule l'emplacement d'une
donnée à partir de sa clé, et qui atteint un coût constant en moyenne. C'est
le hachage : fonctions de hachage, collisions, facteur de charge.
\end{prereqbox}

% =====================================================================
\section*{Références}
\addcontentsline{toc}{section}{Références}
% =====================================================================

\begin{itemize}[nosep]
    \item R. Sedgewick, K. Wayne, \emph{Algorithms}, 4\textsuperscript{e} éd.
          --- §1.4 (Analysis of Algorithms).
    \item T. Cormen, C. Leiserson, R. Rivest, C. Stein,
          \emph{Introduction to Algorithms}, 4\textsuperscript{e} éd.
          --- Ch.~3 (Characterizing Running Times), Ch.~4
          (Divide-and-Conquer).
    \item S. Skiena, \emph{The Algorithm Design Manual},
          3\textsuperscript{e} éd. --- Ch.~2 (Algorithm Analysis).
\end{itemize}

\end{document}
